\documentclass[aspectratio=169]{beamer}
\usepackage{tikz}
\usepackage{graphicx}
\usepackage{listings}
\usepackage{xcolor}
\usepackage{hyperref}
\usepackage[spanish,es-noquoting]{babel}
\usepackage[round]{natbib}

\lstset{
  basicstyle=\ttfamily\small,
  keywordstyle=\color{blue},
  stringstyle=\color{red},
  commentstyle=\color{gray},
  showstringspaces=false
}
% Minimal look
\usetheme{Warsaw}
%\setbeamertemplate{headline}
\setbeamertemplate{footline}{
  \leavevmode%
  \hbox{%
    % Left: Section name
    \begin{beamercolorbox}[wd=.33\paperwidth,ht=2.25ex,dp=1ex,left]{author in head/foot}%
      \hspace*{2ex}\insertsection
    \end{beamercolorbox}%
    % Center: Presentation title
    \begin{beamercolorbox}[wd=.34\paperwidth,ht=2.25ex,dp=1ex,center]{title in head/foot}%
      \insertshorttitle
    \end{beamercolorbox}%
    % Right: Slide numbers
    \begin{beamercolorbox}[wd=.33\paperwidth,ht=2.25ex,dp=1ex,right]{date in head/foot}%
      \insertframenumber{} / \inserttotalframenumber\hspace*{2ex}
    \end{beamercolorbox}%
  }%
  \vskip0pt%
}

\usecolortheme{default}
\usefonttheme{serif}
\graphicspath{{../}{./}}
\setbeamertemplate{caption}[numbered]

% Title info
\title[Handout - Ejercicios Prácticos]{Handout - Ejercicios Prácticos (Módulo 2.4)}
\subtitle{}
\author{Ricardo Reyes-Jimenez}
\institute{Universidad de La Sabana}
\date{\today}

% Background image for title slide
\usebackgroundtemplate{
  \tikz\node[opacity=0.3]{\includegraphics[width=\paperwidth,height=\paperheight]{Images/UniSabana.jpg}};
}

\begin{document}

% Title slide
\begin{frame}
  \titlepage
\end{frame}

% Reset background for normal slides
\usebackgroundtemplate{}



\begin{frame}[fragile]{Ejercicio 1: Los 5 Pilares (Civil - Revisión de Contrato)}


\textbf{Contexto del Caso (Hechos para tu Prompt):}
Tu cliente, ``Inversiones ABC'', firmó un contrato civil de compraventa de un lote comercial con ``Constructora XYZ''. El contrato establecía el pago en 3 cuotas. La segunda cuota se retrasó 5 días por un problema bancario. La Constructora XYZ ha enviado una notificación de resolución automática del contrato basándose en la Cláusula 8.

\textbf{Cláusula 8 del Contrato:} ``El retraso superior a 3 días calendario en el pago de cualquier cuota dará derecho al vendedor a resolver de plano este contrato sin necesidad de requerimiento judicial previo.''

\textbf{Tu Misión:}
Redactar un prompt utilizando los 5 Pilares (Rol, Contexto, Instrucción, Restricciones, Formato) para que la IA evalúe si la cláusula 8 es aplicable o si constituye abuso del derecho según la jurisprudencia civil colombiana.

\end{frame}



\begin{frame}[fragile]{Ejercicio 2: 5 Pilares + CoT + Few-Shot + CARE (Laboral)}


\textbf{Contexto del Caso (Hechos para tu Prompt):}
El empleado ``Juan Pérez'', que lleva 3 años en la empresa, fue despedido alegando ``justa causa'' debido a ausencias reiteradas. Sin embargo, Juan afirma que las ausencias fueron por calamidad doméstica comprobada, pero el departamento de RRHH nunca revisó sus soportes médicos y no se hizo un proceso disciplinario previo.

\textbf{Tu Misión:}
Usando el esquema CARE (Contexto, Acción, Resultado, Evaluación), redacta un prompt que obligue a la IA a realizar un análisis secuencial (Chain-of-Thought) antes de dar el veredicto sobre si el despido fue legal. Además, incluye al menos un ejemplo (Few-Shot) de cómo quieres que redacte la respuesta final (ej. ``En nuestra firma, siempre concluimos con un subtítulo de Recomendación Estratégica: ...'').

\end{frame}



\begin{frame}[fragile]{Ejercicio 3: 5 Pilares + CoT + ADP + Red Team (Contratación)}


\textbf{Contexto del Caso (Hechos para tu Prompt):}
Representas al ``Proveedor P'' (una PYME de tecnología). Le están pidiendo firmar un contrato de servicios con ``Megacorporación S.A.''.

\textbf{Cláusula 12 (Indemnidad):} ``El Proveedor se obliga a mantener indemne a Megacorporación S.A. por cualquier daño directo, indirecto, lucro cesante o daño reputacional derivado de cualquier interrupción del servicio de software, sin límite de cuantía.''

\textbf{Tu Misión:}
1. \textbf{Fase 1 (ADP):} Usa el framework ADP (Actor, Documento, Propósito) y CoT para auditar la Cláusula 12, indicando por qué es desproporcionada y proponiendo una nueva redacción con un límite de responsabilidad (cap).
2. \textbf{Fase 2 (Red Teaming):} Una vez que la IA proponga la nueva redacción mitigada, envíale un segundo prompt ordenándole actuar como el abogado de Megacorporación S.A., pidiéndole que ataque tu nueva redacción y encuentre 3 razones para rechazarla.

\end{frame}



\begin{frame}[fragile]{Ejercicio 4 (Para llevar): Framework TAM (Comunicación / Comercial)}


\textbf{Contexto del Caso (Hechos para tu Prompt):}
Tu firma de abogados debe enviar un requerimiento de pago a ``Distribuidora El Sol'', un cliente histórico de tu representado ``Fábrica La Luna''. La deuda es de \$50,000,000 COP y tiene 60 días de mora. El objetivo es cobrar el dinero rápidamente pero \textbf{sin destruir la relación comercial de 10 años}.

\textbf{Tu Misión:}
Aplica el framework TAM (Tarea, Audiencia, Método) dentro de un prompt para que la IA redacte el correo electrónico perfecto.
\begin{itemize}
  \item \textbf{Tarea:} Solicitar el pago de la deuda.
  \item \textbf{Audiencia:} El Gerente General de ``Distribuidora El Sol'' (relación amistosa pero profesional).
  \item \textbf{Método:} Correo electrónico persuasivo, conciliador pero firme, proponiendo un acuerdo de pago.
\end{itemize}
\end{frame}




\end{document}